% !PRINTED 614
\weberpage{614}
%<<CONT

5. \gesperrt{Zu jedem ganzen Functional des K\"orpers $\Omega$ geh\"ort eine und nur eine Reihe von quadratischen Irrationalzahlen der Discriminante $\Delta$, und die Zahlen dieser Reihe unterscheiden sich um beliebige ganze rationale Zahlen und sind folglich unter einander \"aquivalent. Associirte Formen f\"uhren auf dieselbe Reihe.}

Wenn $\varphi_1$, $\varphi_2$ zwei ganze Functionale in $\Omega$ bedeuten, und
\[
\varphi_1 \varphi_2 = \varphi_3
\]
gesetzt wird, so k\"onnen wir zu jeder dieser drei Formen eine quadratische Irrationalzahl $\omega_1$, $\omega_2$, $\omega_3$ bestimmen, und $\omega_3$ ist durch $\omega_1$, $\omega_2$ bis auf eine additive ganze rationale Zahl bestimmt. Man nennt dann $\omega_3$ aus $\omega_1$ und $\omega_2$ \gesperrt{componirt} oder zusammengesetzt, und setzt in symbolischer Bezeichnung
\weq{3}{\omega_1 \omega_2 = \omega_3.}

Ist $\psi_1$ \"aquivalent mit $\varphi_1$, $\psi_2$ mit $\varphi_2$, so ist nach \S.\,152 auch $\psi_1 \psi_2 = \psi_3$ mit $\varphi_3$ \"aquivalent.

Wenn nun $\eta_1$, $\eta_2$, $\eta_3$ die durch den Satz 5. zu den Formen $\psi_1$, $\psi_2$, $\psi_3$ geh\"origen Irrationalzahlen sind, so bilden nach 3., \S.\,167 die Zahlen $\eta_1$, $\omega_1$; $\eta_2$, $\omega_2$; $\eta_3$, $\omega_3$ drei Paare \"aquivalenter Zahlen und wir k\"onnen daher folgenden Satz aussprechen:

6. \gesperrt{Ersetzt man in einer Composition von quadratischen Irrationalzahlen jedes Element durch eine \"aquivalente Zahl, so geht auch die componirte Zahl in eine \"aquivalente \"uber.}

Dadurch sind die Classen der quadratischen Irrationalzahlen einer Discriminante $\Delta$ zu einer mit der Gruppe der Idealclassen isomorphen \gesperrt{Abel'schen} Gruppe verbunden.

Um die aus zwei gegebenen Zahlen componirte Zahl wirklich zu finden, hat man so zu verfahren.

Wir wollen die nach (1) gebildeten Basen der Formen $\varphi_1$, $\varphi_2$, $\varphi_3$ so bezeichnen:
\[
\begin{array}{lll}
\varphi_1) & \alpha_1 = a_{1,1} & \alpha_2 = a_{1,2} + a_{2,2}\,\Theta \\
\varphi_2) & \beta_1 = b_{1,1} & \beta_2 = b_{1,2} + b_{2,2}\,\Theta \\
\varphi_3) & \gamma_1 = c_{1,1} & \gamma_2 = c_{1,2} + c_{2,2}\,\Theta,
\end{array}
\]
so dass
\[
\omega_1 = \frac{\alpha_2}{\alpha_1}, \quad \omega_2 = \frac{\beta_2}{\beta_1}, \quad \omega_3 = \frac{\gamma_2}{\gamma_1}
\]
%CONT>>
